% -----------------------------------------------------------
\section{Proofs} \label{app:proofs}
% ------------------------------------------------------------

Throughout this appendix, $b_m$ denotes the volume of the unit ball in $\R^m$, $s_{m-1}=mb_m$ the $(m-1)$-dimensional surface measure of the unit sphere $S^{m-1}\subset\R^m$, and $B(\bm{x},r)$ the closed ball of radius $r$ centred at $\bm{x}$.

\subsection{Proof of \texorpdfstring{\cref{prop:order}}{Proposition 1}}

Write $\mathcal{Y}_n=\{\bm{y}_1,\dots,\bm{y}_n\}$ and $k_n=\lceil\gamma(n-1)\rceil$, and recall from \cref{sec:order} that $C_n(r)$ is the number of components of the graph on $\mathcal{Y}_n$ that joins pairs at distance at most $r$.

\begin{proof}[Proof of \cref{prop:order}]
  \emph{Step 1: the component count.} For $t>0$ put $\Xi_n(t):=n^{-1}C_n(tn^{-1/m})$ and
  $$
  G(t):=\int_{\R^m}\kappa\big(t^mf(\bm{x})\big)f(\bm{x})\,d\bm{x}.
  $$
  Joining at distance $r_n=tn^{-1/m}$ keeps $nr_n^m=t^m$ fixed, which is the thermodynamic regime, so the law of large numbers for component counts, \citet[Theorem 3.13, an application of their Theorem 2.1]{MitschePenrose2021}, gives $\Xi_n(t)\to G(t)$ in probability, for every $t>0$ and every density $f$. Their limit function for component counts is $\kappa$, and their convergence is in fact complete and in $L^2$. The same conclusion follows from the law of large numbers for the empirical distribution of the rescaled edge lengths of $T$ \citep[Theorem 2.3(i), first obtained by \citealp{Penrose1996}]{PenroseYukich2003}, since $n-C_n(r)$ is the number of edges of $T$ of length at most $r$; the component-count form is used here because its limit is expressed through $\kappa$.

  \emph{Step 2: properties of $G$.} The origin forms a component on its own exactly when $\mathcal{P}_\lambda$ has no point in $B(\bm{0},1)$, so $\kappa(\lambda)\ge\P(|\mathcal{C}_\lambda|=1)=\exp(-b_m\lambda)$. Its component contains the $N\sim\text{Poisson}(b_m\lambda)$ points of $\mathcal{P}_\lambda$ in $B(\bm{0},1)$, so $\kappa(\lambda)\le\E[1/(1+N)]=(1-\exp(-b_m\lambda))/(b_m\lambda)$. Integrating,
  $$
  G_-(t):=\int_{\R^m}\exp(-b_mt^mf(\bm{x}))f(\bm{x})\,d\bm{x}\ \le\ G(t)\ \le\ G_+(t):=\frac{1}{b_mt^m}\int_{\R^m}\big(1-\exp(-b_mt^mf(\bm{x}))\big)\,d\bm{x}.
  $$
  By dominated convergence $G_-$ and $G_+$ are continuous and strictly decreasing, tending to $1$ as $t\to0$ and to $0$ as $t\to\infty$; for $G_+$ this uses that $(1-\exp(-u))/u$ decreases strictly from $1$ to $0$ on $(0,\infty)$. Hence $t_-$ and $t_+$ are well defined, and $G(t)\to1$ as $t\to0$ and $G(t)\to0$ as $t\to\infty$.

  The function $\kappa$ is continuous on $(0,\infty)$, because $\lambda\mapsto\lambda\kappa(\lambda)$ is Lipschitz \citep[Proposition 2.3(ii) and Section 3.6]{MitschePenrose2021}, and dominated convergence then makes $G$ continuous. To see that $G$ is strictly decreasing, fix $\tau>0$, take one realisation of $\mathcal{P}_\tau\cup\{\bm{0}\}$, and let $\mathcal{C}(s)$ be the component of the origin when points are joined at distance $s$, so that $\kappa(\tau s^m)=\E[1/|\mathcal{C}(s)|]$. Components only grow with $s$, and when $s_1<s_2$ and the nearest point to the origin lies at a distance in $(s_1,s_2]$, the origin is alone at scale $s_1$ and not at scale $s_2$, so
  $$
  \kappa(\tau s_1^m)-\kappa(\tau s_2^m)\ \ge\ \tfrac12\big(\exp(-b_m\tau s_1^m)-\exp(-b_m\tau s_2^m)\big)>0.
  $$
  Hence $G$ is strictly decreasing, because its integrand decreases strictly in $t$ wherever $f>0$, and \eqref{eq:qgamma} has a unique solution $q_\gamma(f)\in(0,\infty)$. Since $G_-(q_\gamma(f))\le G(q_\gamma(f))=1-\gamma=G_-(t_-)$ and $G_-$ is decreasing, $q_\gamma(f)\ge t_-$, and the same argument with $G_+$ gives $q_\gamma(f)\le t_+$.

  \emph{Step 3: the quantile.} Write $q=q_\gamma(f)$ and fix $\varepsilon\in(0,q)$, so that $G(q-\varepsilon)>1-\gamma>G(q+\varepsilon)$. Since $h_n\le r$ exactly when $C_n(r)\le n-k_n$, and $\gamma(n-1)\le k_n<\gamma(n-1)+1$,
  $$
  \P(n^{1/m}h_n\le q-\varepsilon)\le\P\big(\Xi_n(q-\varepsilon)\le1-\gamma+\gamma/n\big)
  \quad\text{and}\quad
  \P(n^{1/m}h_n>q+\varepsilon)\le\P\big(\Xi_n(q+\varepsilon)>1-\gamma-1/n\big).
  $$
  Both tend to zero by Step 1, so $n^{1/m}h_n\to q$ in probability.

  Finally, for a locally finite $\mathcal{X}\subset\R^m$ and $\bm{x}\in\mathcal{X}$, let $\psi_t(\bm{x};\mathcal{X})=\mathbf{1}\{\mathcal{X}\cap B(\bm{x},t)\neq\{\bm{x}\}\}$ indicate that some other point of $\mathcal{X}$ lies within distance $t$ of $\bm{x}$. It is bounded, unchanged when $\bm{x}$ and $\mathcal{X}$ are translated together, and depends only on the points within distance $t$ of $\bm{x}$, so \citet[Theorem 2.1]{PenroseYukich2003} applies and gives
  $$
  \frac1n\sum_{i=1}^n\psi_t(n^{1/m}\bm{y}_i;n^{1/m}\mathcal{Y}_n)\to\int_{\R^m}\P\big(\mathcal{P}_{f(\bm{x})}\cap B(\bm{0},t)\neq\emptyset\big)f(\bm{x})\,d\bm{x}=1-G_-(t)
  $$
  in probability. The left side is the proportion of observations whose nearest neighbour lies within $tn^{-1/m}$, so the argument of Step 3, with $G_-$ in place of $G$, gives $n^{1/m}D_{(k_n)}\to t_-$ in probability, as stated in \cref{sec:order}.
\end{proof}

\subsection{Proof of \texorpdfstring{\cref{cex:fattails}}{Counterexample 1}}

Let $f(\bm{x})=c_m/(1+\|\bm{x}\|^{m+1})$ on $\R^m$ with $m\ge1$, let $\bm{Y}$ have density $f$, and write $a_m=c_ms_{m-1}$. We use two elementary bounds. First, $\|\bm{Y}\|$ has density $a_mr^{m-1}/(1+r^{m+1})$, which for $r\ge1$ lies between $a_m/(2r^2)$ and $a_m/r^2$, so
\begin{equation}\label{eq:tailprob}
  \frac{a_m}{2R}\le p(R):=\P(\|\bm{Y}\|\ge R)\le\frac{a_m}{R},\qquad R\ge1.
\end{equation}
Second, $f(\bm{z})\le c_m\|\bm{z}\|^{-(m+1)}$, so for every $\bm{x}$ and every $s\in(0,\|\bm{x}\|)$,
\begin{equation}\label{eq:ballprob}
  \P\big(\|\bm{Y}-\bm{x}\|\le s\big)=\int_{B(\bm{x},s)}f(\bm{z})\,d\bm{z}\le\frac{b_mc_ms^m}{(\|\bm{x}\|-s)^{m+1}}.
\end{equation}

\begin{proof}[Proof of \cref{cex:fattails}]
  Put $R_n=a_mn^{3/4}$, $s_n=n^{1/8}$ and $K_n=n-\omega_n=\lfloor n^{1/4}/4\rfloor$. Let $N_n$ be the number of observations with $\|\bm{y}_i\|\ge R_n$, and let $Y_n$ be the number of pairs $\{i,j\}$ with $\|\bm{y}_i\|\ge R_n-s_n$, $\|\bm{y}_j\|\ge R_n-s_n$ and $\|\bm{y}_i-\bm{y}_j\|\le s_n$. Throughout, $n$ is large enough that $R_n-2s_n\ge\max(1,R_n/2)$.

  \emph{Step 1: the number of distant observations.} $N_n$ is binomial with $n$ trials and success probability $p(R_n)$, which by \eqref{eq:tailprob} lies in $[n^{-3/4}/2,\,n^{-3/4}]$. Its mean $\mu_n$ therefore satisfies $n^{1/4}/2\le\mu_n\le n^{1/4}$, its variance is at most $\mu_n$, and $K_n\le\mu_n/2$. By Chebyshev's inequality,
  $$
  \P(N_n\le K_n)\le\P\big(|N_n-\mu_n|\ge\mu_n/2\big)\le\frac{4}{\mu_n}\le8n^{-1/4}\to0.
  $$

  \emph{Step 2: no close pairs among distant observations.} Conditioning on the position of the first point of a pair and applying \eqref{eq:tailprob} and \eqref{eq:ballprob},
  $$
  \E[Y_n]\le\binom{n}{2}\,p(R_n-s_n)\sup_{\|\bm{x}\|\ge R_n-s_n}\P\big(\|\bm{Y}-\bm{x}\|\le s_n\big)
  \le\frac{n^2}{2}\cdot\frac{a_m}{R_n-s_n}\cdot\frac{b_mc_ms_n^m}{(R_n-2s_n)^{m+1}}
  \le\frac{2^{m+1}a_mb_mc_m\,n^2s_n^m}{R_n^{m+2}}.
  $$
  The right side is $C_m\,n^{2+m/8-3(m+2)/4}=C_m\,n^{1/2-5m/8}$ with $C_m$ depending only on $m$, and tends to zero for every $m\ge1$. Hence $\P(Y_n\ge1)\le\E[Y_n]\to0$.

  \emph{Step 3: conclusion.} Suppose $N_n>K_n$ and $Y_n=0$. Any observation within distance $s_n$ of an observation $\bm{y}_i$ with $\|\bm{y}_i\|\ge R_n$ has norm at least $R_n-s_n$, so the pair would be counted in $Y_n$, and $Y_n=0$. Each of the $N_n$ observations beyond $R_n$ is therefore a component on its own when observations are joined at distance $s_n$, so $C_n(s_n)\ge N_n>K_n=n-\omega_n$. By \cref{sec:order}, exactly $n-C_n(s_n)$ edges of $T$ have length at most $s_n$, which is fewer than $\omega_n$, so $e_{(\omega_n)}>s_n$. Therefore
  \begin{equation*}
  \P\big(e_{(\omega_n)}>n^{1/8}\big)\ge1-\P(N_n\le K_n)-\P(Y_n\ge1)\to1.\qedhere
  \end{equation*}
\end{proof}
